\section{Anytime-valid risk certification}\label{sec:sequential}
\subsection{Maturation and the observable confidence sequence}
An event enters inference once, when its age reaches $H$, even if its label was resolved earlier. Its category is then permanent. This avoids selecting earlier statistical observations merely because they resolve faster. Let $Y_1,Y_2,\ldots$ denote the categories in this maturity order, with filtration $\mathcal F_n$. Assume that there is one fixed categorical law $p$ such that
\begin{equation}\label{eq:stable}
 \Ebb[\ind\{Y_n=j\}\mid\mathcal F_{n-1}]=p_j
 \quad\text{for every category }j\text{ and every }n.
\end{equation}
IID categorical draws meet this condition. Any certificate pertains to a stream governed by a known fixed policy / checking protocol / time-horizon / meaning of labels. In particular, Eq.~\eqref{eq:stable} excludes any general data-drift and general informative data-permutations.

Similarly, $\theta$ and $u$ are tied to the law of all (matured) data points, as fixed by Eq.~\eqref{eq:stable}. It further states an assumption about the distribution of the \emph{categories} of observed data being fixed, enabling monitoring sequentially over time. Note this would not imply the true risk parameter is unchanged.

Let $S_{j,n} = \sum_{i \leq n} \ind\{Y_i = j\}$. Then the Jeffreys-prior LR by category for $x \in (0,1)$ is given by
\begin{equation}\label{eq:mixture}
 M_{j,n}(x)=\frac{B(S_{j,n}+1/2,n-S_{j,n}+1/2)}
 {B(1/2,1/2)x^{S_{j,n}}(1-x)^{n-S_{j,n}}},
\end{equation}
where $B$ is the beta function. Define raw intervals as the closures of $\{x:M_{j,n}(x)<d/\delta\}$, take their running intersections over all prefixes, and intersect their product with the probability simplex:
\begin{equation}\label{eq:CS}
 \Cset_n=\{p\in\mathbb R_+^d:\sum_jp_j=1,
              \ \ell_{j,n}\le p_j\le v_{j,n}\text{ for every }j\}.
\end{equation}
The construction allocates $\delta/d$ per category and gives
\begin{equation}\label{eq:CSguarantee}
 \Pbb\{p\in\Cset_n\text{ for every }n\}\ge1-\delta.
\end{equation}
\citet{howard2021} and \citet{ramdas2023} outline these mixture/martingale methods and maximal inequalities for confidence sequences respectively. We adopt a coordinate-wise version since there is a convenient covering set for the collection of (coordinate-wise) intervals; cf. \citep{lindon2022} for joint confidence regions over the full multinomial distribution. For endpoint null hypotheses ($p_j\in\{0,1\}$), we use the appropriate limit of the likelihood. Root-finding is done numerically, giving us upper and lower bracketing values. In our accelerated code, this starts from a standard FP root sign determination which defers to Arb if no sign is found, so valid bracketed bounds require correct FP signs + bounds on Arb output. Confidence region is refreshed after every maturity event (although optimisation of the risk is typically done only at specific times).

\subsection{Conservative projection and computation}
The ideal upper-risk projection is
\begin{equation}\label{eq:idealprojection}
 U_n^{\rm exact}(\rho)=\sup_{p'\in\Cset_n}\sup\Thset(p',\rho),
\end{equation}
where incompatible $p'$ contribute no risks. An empty feasible union produces no certificate. This is a confidence projection of the identified set, not a plug-in estimate of its upper endpoint.

The implementation first encloses all observable summaries from $\Cset_n$ in an aggregate envelope. Direct category sums and complementary sums give simultaneous intervals for $A,G,c$. In particular,
\begin{equation}\label{eq:aggregate}
 A_L=\max\Big\{\sum_k\ell_{a_k},\ 1-\sum_kv_{b_k}-v_c,\ 0\Big\},\quad
 A_U=\min\Big\{\sum_kv_{a_k},\ 1-\sum_k\ell_{b_k}-\ell_c,\ 1\Big\},
\end{equation}
with analogous formulas for $G$ and $c$. By coordinatewise monotonicity of $q(a,b)$,
\begin{align}
 T_L&=\sum_k q(\ell_{a_k},\ell_{b_k}),\label{eq:entropyenv}\\
 T_U&=\min\left\{\sum_kq(v_{a_k},v_{b_k}),\
 \min\big(1,\sum_k(v_{a_k}+v_{b_k})\big)\log2,\ \log2\right\}.\nonumber
\end{align}
Thus $\Eset_n=\{(A,G,T):A\in[A_L,A_U],G\in[G_L,G_U],1-A-G\in[c_L,c_U],T\in[T_L,T_U]\}$ contains every true summary on the event in \eqref{eq:CSguarantee}. It may contain combinations unattainable by a categorical law. Maximizing over this envelope is a deliberate relaxation.

Because the information constraint decreases in $T$, the relaxed upper projection sets $T=T_U$ and searches over $(A,G,u)$:
\begin{equation}\label{eq:relaxedprojection}
 V_n(\rho)=\sup\left\{A+u:
 \begin{array}{l} A\in[A_L,A_U],\ G\in[G_L,G_U],\ c=1-A-G\in[c_L,c_U],\\
 0\le u\le c,\ h(A+u)-T_U-q(u,c-u)\le\rho
 \end{array}\right\}.
\end{equation}
Any nonempty feasible union appearing in \eqref{eq:idealprojection} is included in the relaxation considered here. We use numerical search to obtain an upper enclosure to $V_n$, and invoke the anytime (time-uniform) property whenever such enclosures are guaranteed to be conservative, per Theorem~\ref{thm:anytime}.

The deterministic branch-and-bound operates over boxes parameterised by the three sufficient statistics. If a box's feasibility conditions fail, or its enclosed minimum information content is >$\rho$, we eliminate it from further exploration. All un-eliminated boxes then add their corresponding upper objective value $\min(1,A_{\rm hi}+u_{\rm hi})$ to the overall upper bound. Joint tables potentially meeting feasibility requirements may provide an incumbent lower objective value but cannot be used to certify the overall bound. Information-content bounds are obtained using Arb interval/ball arithmetic, with 128-bit Arb for projection \citep{johansson2017}, and standard FP operations (residual-corrected summation, outward-rounded summaries, singletons) in the code at present, except to perform exact checks on candidate joint tables. We outline the assumptions required for individual pieces in order to preserve this bounding guarantee, and summarise the fast-filter for likely infeasible boxes in the Supplement, although we were unable to establish a generic bound on the errors introduced by that approximation from inspection of the code. Therefore Theorem~\ref{thm:anytime} relies conservatively both on the ability to project intervals and obtain conservative enclosing inversions. Empirical testing and observation of rates of certified solutions is not sufficient in lieu of such guarantees. Furthermore, while box size will tend to constrain the total possible change in objective value, a tight box does not necessarily entail a low bound-objective gap because we have no analogous guarantee about local variation of the total information content bound.

We thus output the maximum upper bound upon any of termination conditions (full tree explored, stagnation, minimum box size, numerical caution). We will indicate when the solution has converged based on our chosen upper-incumbent gap threshold. To solve for decisions, the search may terminate when the upper objective is $\leq \beta$, or a feasible-solution upper objective estimate is >$\beta$. In the latter case, assuming that solution itself can be validated as numerically feasible, we can refrain from claiming a certified bound but cannot infer that the true objective is >$\beta$. We defer details on this and proof of the bound-invariant to the Supplement.

\begin{theorem}[Time-uniform sensitivity certificate]\label{thm:anytime}
Suppose \eqref{eq:stable} holds, the true full law satisfies \eqref{eq:pis}, the computed confidence regions contain the exact regions in \eqref{eq:CS}, and each reported numerical bound conservatively encloses the projection described above. For a fixed analysis partition and its per-category allocation,
\begin{equation}\label{eq:anytime}
 \Pbb\{\theta\le\widehat U_n(\rho)\text{ at every valid reported update }n\}\ge1-\delta.
\end{equation}
Consequently, for any monitoring rule that issues a certificate only when $\widehat U_n(\rho)\le\beta$, the probability of ever certifying when $\theta>\beta$ is at most $\delta$.
\end{theorem}
\begin{proof}
For any genuine $p_j\in(0,1)$, \eqref{eq:mixture} is a nonnegative mixture martingale with unit initial value under \eqref{eq:stable} (or supermartingale if $p_j\in\{0,1\}$). Applying the maximal inequality and a union bound yields \eqref{eq:CSguarantee}. Conditional on that simultaneous event, the actual summary and the true terminal completion are valid candidates in all projections. Moreover, data-processing ensures that the coarse information is bounded above by $\rho$, so the corresponding projection/empirical bounds are each larger than $\theta$ throughout. Hence spurious certificates can only arise by escaping that simultaneous event.
\end{proof}

The statement is per prespecified partition. Comparing four partitions in a simulation does not allocate a joint four-partition error budget. Choosing a partition from the same data and issuing whichever certificate is most favorable requires a simultaneous allocation across those procedures or a common finest-law region. Likewise, the theorem concerns monitoring within one stable stream, not a familywise guarantee over many streams.

\subsection{Decision states}
Given sufficient data and valid numerics, the gate then requires at least 200 matured and at least 50 resolved counts. The subsequent decisions are made sequentially: declare \emph{model incompatible}, if a conservative lower bound to all admissible compatibility thresholds is larger than $\rho$; else declare \emph{intrinsically uncertifiable} if a conservative lower bound to all compatibility-admissible information-limited risks is larger than $\beta$; else declare \emph{certified} if $\widehat U_n\leq\beta$, and \emph{uncertified} else. Insufficient count will trigger \emph{insufficient evidence}. The strong intrinsic state is withheld if zero resolved population mass remains plausible. Lower bounds must be computed conservatively to support the first two findings (the comparison protections by themselves do not suffice).

Confidence region becoming empty corresponds to an error event for our procedure based on confidence sequences. Our procedure ceases to generate output but remains part of the running number of simulations (not restarted with a new random seed). Note the difference to e.g. a computational error. None of the states indicate validation of the unknown sensitivity under model compatibility.
